% ---------------------------------------------------------------------
\subsection{One sounding pipe (\fn{OrganEngine::renderRange})}
\label{sec:math-voice}
A \emph{voice} is one pipe that sounds: it reads the pipe's mono sample, shapes it sample by sample, and adds it to
the output with gains that put it in its place. This subsection follows one voice through one block in the order of
the code. Many small filters below are the same one-pole low-pass, so it is defined once:
\begin{equation}
  y_n = y_{n-1} + a\,(x_n - y_{n-1}),\qquad a = 1 - e^{-2\pi f_c/f_s}.
  \label{eq:onepole}
\end{equation}
This is the impulse-invariant discretisation of an $RC$ low-pass with corner $f_c$; its step response reaches
$1-1/e$ after $f_s/(2\pi f_c)$ samples. The same update with $a = 1-e^{-T/\tau}$, applied once every $T$ seconds,
is an exponential glide with time constant $\tau$; the code uses it wherever a value must move smoothly instead of
jumping (for example $a=0.03$ every 16 samples is $\tau = -16/(f_s\ln 0.97) \approx \SI{11}{ms}$ at \SI{48}{kHz}).

\paragraph{Starting a voice (\fn{startVoice}).} A free voice is taken; if all are busy, the oldest voice whose key is
already released is stolen, and only if every voice is held, the oldest of all. The stolen voice is not cut: it
fades out linearly over \SI{5}{ms} while the new note starts. The voice is given a fixed random detune
$\rho\in[-1,1]$ seeded by stop and key (so the same pipe is always out of tune the same way), the key position
$\kappa = 2k/40-1\in[-1,1]$ as a modulation source, and its loop crossfade length
\begin{equation}
  X = \operatorname{clamp}\bigl(\min(1440,\ L/3),\ 16,\ \ell_s\bigr)\ \text{samples of the file},
\end{equation}
with the loop $[\ell_s,\ell_e)$ of length $L = \ell_e-\ell_s$: at most \SI{30}{ms} at the file rate of
\SI{48}{kHz}, at most a third of the loop, and never longer than the part before the loop.

\paragraph{The reading rate.} The playback rate in samples of the file per output sample is
\begin{equation}
  r_n = \underbrace{\frac{f_\mathrm{file}}{f_s}\,\rho_\mathrm{native}\,\rho_\mathrm{cal}}_{r_0}
        \cdot 2^{\,c_n/1200}\cdot D_n ,
\end{equation}
where $f_\mathrm{file}/f_s$ converts the file's \SI{48}{kHz} to the host rate, $\rho_\mathrm{native}$ is the pipe's
static transposition (0.5 for the Praestant~16′, which reads the 8′ samples an octave lower, and the shift of a
substituted pipe), $\rho_\mathrm{cal}$ the optional correction of the recorded pitch to the nominal one
(\fn{samplePlaybackRate}; 1 when it is off), $D_n$ the Doppler factor and $c_n$ the sum of all pitch offsets in cents,
re-evaluated every 16 samples:
\begin{multline}
  c = \underbrace{\Delta_k + c_S + \delta\rho}_{\text{tuning}}
     \;+\; \underbrace{c_\ell + c_\mathrm{pure} + 38\,w_r^2\,\sigma_w(0.4+0.6|\rho|)}_{\text{length, pure tuning, wear}}
     \;+\; \underbrace{c_\mathrm{wind} + c_\mathrm{wob} + c_\mathrm{gust} + c_\mathrm{trem}}_{\text{wind}}\\
     \;+\; \underbrace{c_\mathrm{mod}(n) + c_\mathrm{MPE}}_{\text{modulation}} .
\end{multline}
$\Delta_k$ is the key's tuning shift (\S\ref{sec:math-tuning}), $c_S$ the \ui{Stimmton}, $\delta$ the
\ui{Verstimmung} (each pipe detuned by its own $\delta\rho$). $c_\ell = -1200\log_2\lambda$ follows the pipe's length
factor $\lambda$ (\ui{Wachsen \& Brechen}: a pipe twice as long sounds an octave lower; it glides with
$\tau=\SI{30}{ms}$, so a pipe that breaks while sounding jumps without a click). $c_\mathrm{pure}$ comes from
\ui{Reinheit} and the wear term from \ui{Verschleiß} of the living organ (\S\ref{sec:live}); $\sigma_w=-1$ for
60\,\% of the pipes, because worn pipes mostly go flat. The modulation term is interpolated linearly across the
block from its value at the start to the value at the end, $c_\mathrm{mod}(n) = c_0 + (c_1-c_0)\,n/N$, so a pitch LFO
produces a smooth glide instead of a staircase of 64-sample steps. The MPE bend glides with
$a=0.2$ per 16 samples ($\tau\approx\SI{1.5}{ms}$).

\paragraph{The wind terms.} With the stop wind $w$ (the stop's \ui{Wind} slider plus its modulation, times the
bellows pressure and the pallet opening of the living organ, smoothed with $\tau\approx\SI{11}{ms}$) and the lack of
wind $\ell = 1-w$, the classic model of \S\ref{sec:stopwind} gives
$c_\mathrm{wind} = -75\,\ell^{3/2} + 35\,\ell^2 b(t)$ and the level factor $a = w^{1.6}(1+0.4\,\ell\,b(t))$, with the
breathing $b(t)$ a \fn{Wobble} (below) advanced at $1+3\ell$ times its own speed. The exponent $1.6$ makes the level
fall faster than linearly: a half-drawn stop is about $\SI{-10}{dB}$, not $\SI{-6}{dB}$. Then:
\begin{itemize}[nosep]
  \item \emph{Windschwankung} with the amount $A = \alpha\omega + 0.45w_r^2$ (the degradation amount times the
    wobble, plus wear): $c_\mathrm{wob} = (30A + 320A^4)\,u(t)$, level $\times\max(0, 1 + (0.3A+0.6A^4)u(t))$, speed
    $1+5A^3$. The quartic term is negligible up to about half and dominates near the top, which is why the slider
    gets disproportionately wilder towards 100\,\%: at $A=0.5$ the pitch swings $\pm\SI{35}{ct}$, at $A=1$ by
    $\pm\SI{350}{ct}$.
  \item \emph{Gusts} of the global wind above 35\,\%: $\tau_g = \bigl(\max(0,(W-0.35)/0.65)\bigr)^2$,
    $c_\mathrm{gust} = 190\,\tau_g\,\nu(t)$ and level $\times\max(0.05,\,1+0.75\tau_g\nu(t))$, where $\nu$ is one
    \fn{smoothNoise} (\S\ref{sec:math-motion}, seed 4242) run at $0.5+6\tau_g$ times its speed. One curve is
    computed for all pipes (one value per 16 samples), because the pipes share one wind chest and must shake together.
  \item \emph{Tremulant}: the wave $q(\theta) = \bigl(\sin\theta + 0.25\sin(2\theta+0.7)\bigr)/1.2$ with
    $\theta = 2\pi f_\mathrm{trem} t$, shared by all pipes; with the depth $D$ (the global one plus the stop's own)
    $c_\mathrm{trem} = 16D\,q$ and level $\times(1+0.32D\,q)$. The second harmonic with its phase offset makes the
    wave asymmetric, as the beat of a real tremulant valve is.
\end{itemize}
The wind also closes a one-pole low-pass (\ref{eq:onepole}) with $f_c = 250 + 16000\,w^2$\,Hz (bypassed above
$w=0.995$), so a starving pipe gets duller.

\paragraph{Wobble (\fn{Wobble}).} The slow random movements of the wind are sums of three sines with random rates
$f_1\in[0.15,0.55)$, $f_2\in[0.7,1.9)$, $f_3\in[2.5,5.5)$\,Hz and random phases, drawn once per voice:
\begin{equation}
  u(t) = \bigl(\sin 2\pi\Phi_1 + 0.5\sin 2\pi\Phi_2 + 0.15\sin 2\pi\Phi_3\bigr)/1.65,\qquad
  \Phi_i \mathrel{+}= f_i\,s\,\Delta t,
\end{equation}
where the phases are integrated, so the speed factor $s$ can change without a jump. The weights fall with the
frequency: the movement is mostly slow with a little faster flutter on top, as wind is.

\paragraph{Reading the sample (\fn{readSample}).} The sample is read at the fractional position $p = i + f$ with
four-point cubic Hermite interpolation (Catmull--Rom): the cubic through $x_0, x_1$ whose slopes at the two points
are the central differences $(x_1-x_{-1})/2$ and $(x_2-x_0)/2$,
\begin{gather}
  y = \bigl((c_3 f + c_2) f + c_1\bigr) f + x_0,\\
  c_1 = \tfrac12(x_1 - x_{-1}),\quad c_2 = x_{-1} - \tfrac52 x_0 + 2x_1 - \tfrac12 x_2,\quad
  c_3 = \tfrac12(x_2 - x_{-1}) + \tfrac32(x_0 - x_1).
\end{gather}
It reproduces quadratics exactly, is continuous in value and slope from one interval to the next, and costs four
multiply-adds in Horner form. Its slight attenuation towards the Nyquist frequency is harmless because the pipes are
read at nearly their own rate. There is no anti-alias filter: the only transposition upwards is the tuning shift
(at most a few tens of cents), which moves next to nothing above $f_s/2$.

\paragraph{The loop.} While the voice is held and the position comes within $X$ of the loop end $\ell_e$, the
reader fades linearly into the same place one loop earlier,
\begin{equation}
  x = (1-t)\,s(p) + t\,s(p-L),\qquad t = \frac{p - (\ell_e - X)}{X},
\end{equation}
and at $p\ge\ell_e$ jumps back by $L$, to where the fade has arrived. A linear (equal-gain) fade keeps the amplitude
because the loop points were chosen where the signal is highly correlated with itself (\S\ref{sec:loop}); an
equal-power fade would bulge by up to $\SI{3}{dB}$ for correlated signals.

\paragraph{Release (\fn{releaseVoice}).} When the key goes up, the reader fades over $N_r = 0.04f_s$ samples
from wherever it is into a second reader that starts \SI{40}{ms} before the recorded release point $p_\mathrm{rel}$,
\begin{equation}
  x = (1-g)\,s(p) + g\,s(p_\mathrm{rel}-0.04 f_\mathrm{file} + j\,r),\qquad g = j/N_r,\quad j = 0\dots N_r,
\end{equation}
and then continues from the second reader without the loop, so the note ends with the recorded sound of its pallet
closing. With the glitch \ui{Nulldurchgang} $z$ the fade is shortened to $0.04f_s(1-0.995z)$: a hard cut.

\paragraph{Attack shaping (\ui{Ansprache}).} With the onset $p_a$ and the end of the measured attack $p_b$ of the
pipe, the envelope $e = \tfrac12 + \tfrac12\cos(\pi u)$, $u = (p-p_a)/(p_b-p_a)$ clamped to $[0,1]$, falls from 1 to 0
over the attack. It is a function of the read position, not of time, so it stays on the attack also when the pipe
is transposed or the reader is slowed. The four controls act as
\begin{equation}
  x \leftarrow \bigl(x + 3s\,e\,(x - \mathrm{LP}_{\SI{2.5}{kHz}}\,x)\bigr)\,(1+3g^+e)\,(1-g^-e)\,\bigl(1-q(1-e)\bigr)
\end{equation}
(sharpness $s$, gain $g = g^+-g^-$, \ui{Nur Ansprache} $q$). The repeated chiff is a second reader over
$[p_a, p_a + \ell_c]$, $\ell_c = \min(p_b-p_a,\, 0.09f_\mathrm{file})$, windowed by a Hann window of its length in output
samples $\ell_c/r$, the $i$-th repeat scaled by $0.82^i$, with a random pause of $0.55$--$0.95$ burst lengths between
repeats.

\paragraph{Freeze (\fn{grainSample}).} A frozen voice becomes a granular cloud read from the pipe's sustain (the loop
if there is one, else from \SI{0.1}{s} after the attack to the release). Every $\tfrac13\ell_g\cdot U(0.7,1.3)$
samples a new grain starts at a uniformly random place, with the length $\ell_g = (0.06 + 0.5\,s)f_s$ and its own
rate factor $2^{25 s\,b/1200}$, $b\in[-1,1]$, so up to $\pm25$ cents at full smear $s$. Each grain is windowed by a
Hann window; on average three overlap, and Hann windows at a hop of a third of their length sum to $1.5$, hence the
division by $1.5$. The change from the normal reader to the cloud is a linear crossfade over \SI{150}{ms}.

\paragraph{The filters after the reader.} The signal then passes the wind low-pass and the air low-pass in series,
both of type (\ref{eq:onepole}). The air corner falls with the distance $d$:
\begin{equation}
  f_\mathrm{air} = \operatorname{clamp}_{[1200,\,20000]}\Bigl(\frac{\SI{20}{kHz}}{1 + A\,d/\SI{5}{m}}\Bigr)\cdot 2^{6m_b},
\end{equation}
with the amount $A$ (\ui{Luft}) and the brightness modulation $m_b$ (6 octaves per unit). Real air absorbs high
frequencies exponentially with distance and more the higher they are; a one-pole whose corner falls as $1/d$ is a
cheap stand-in that gets the direction of the effect right. \ui{Luft umkehren} $\iota$ mixes
$y = y_\mathrm{air} + 2\iota\,(y_\mathrm{wind} - y_\mathrm{air})$: at $\iota = \tfrac12$ the highs the air took are
given back exactly, at $\iota = 1$ they are added twice, so far pipes become \emph{brighter} than near ones. The level
factor $a$ of the wind terms is applied here, and, in the unstable wind model, the mouth tone of
\S\ref{sec:windchaos} is added after its own pass through the air filter (it belongs to the pipe and is as far
away).

\paragraph{The upper-partial layer.} A second-order Butterworth high-pass splits off the partials above $f_h$. It is
the state-variable filter in the trapezoidal (topology-preserving) form \citep{zavalishin2012}, which keeps its
analogue response up to $0.45f_s$ and can be retuned per sample without becoming unstable:
\begin{gather}
  g = \tan(\pi f_h/f_s),\quad k = \sqrt2,\quad a_1 = \frac{1}{1+g(g+k)},\quad a_2 = g a_1,\quad a_3 = g a_2,\\
  v_3 = x - s_2,\quad v_1 = a_1 s_1 + a_2 v_3,\quad v_2 = s_2 + a_2 s_1 + a_3 v_3,\\
  s_1 \leftarrow 2v_1 - s_1,\quad s_2 \leftarrow 2v_2 - s_2,\\
  h = x - k v_1 - v_2 .
\end{gather}
$v_1$ is the band-pass and $v_2$ the low-pass output, $h$ the high-pass, and $k=\sqrt2 = 1/Q$ makes it Butterworth.
The output is $x + (G\,t - 1)h$: the layer is replaced by itself times the gain $G$ and the transient factor
$t = \operatorname{clamp}_{[0.2,5]}\bigl((E_f/E_s)^{1.5\sigma}\bigr)$, where $E_f$ and $E_s$ are peak followers of
$|h|$ with attack and release time constants of $0.7$/\SI{6}{ms} and $25$/\SI{120}{ms}. During an attack the fast
follower runs ahead of the slow one, $E_f/E_s>1$, so $\sigma>0$ lifts the beginnings and $\sigma<0$ keeps the
sustained surface.

\paragraph{Release puff (\ui{Abriss}).} When the pallet closes, white noise passes a one-pole low-pass at $2.5f_c$ and
a two-pole resonator
\begin{equation}
  y_n = x_n + 2R\cos\omega\,y_{n-1} - R^2 y_{n-2},\qquad \omega = 2\pi f_c/f_s,\qquad R = e^{-\pi f_c/(1.4 f_s)},
\end{equation}
whose poles at $Re^{\pm i\omega}$ give a peak at $f_c$ with the bandwidth $f_c/1.4$, i.e.\ $Q\approx1.4$. Its gain is
scaled by $1-R$, which normalises the peak gain $\approx 1/(1-R)$ of the resonator, so pipes of all pitches puff
equally loud; the envelope decays by the factor $e^{-1/(\tau f_s)}$ per sample.

\paragraph{Random events: crackle and glitches.} Crackle, stutter and clip bursts arrive as Poisson processes: the
waiting time to the next event is drawn as $T = -\ln(u)/\lambda$ seconds with $u$ uniform, which is exponentially
distributed with mean $1/\lambda$ and makes the events independent of each other (no hidden rhythm). The rates are
$\lambda = $ \ui{Dichte} for crackle, $2.5\,g_s^2$ per second for stutter and $3\,g_c^2$ for clipping (the squares keep
the lower half of the sliders sparse). A crackle is a burst of white noise of amplitude $0.2+0.8u^2$ (mostly small,
sometimes large) decaying with the time constant \ui{Länge}. A stutter repeats a slice of \SIrange{20}{150}{ms}
one to five times by setting the reader back. A clip burst lasts \SIrange{8}{78}{ms} and drives the signal by a factor
$\delta\in[4,\,4+30g_c]$ into a hard limit at $\pm0.02$, with the make-up gain $1.5/\sqrt\delta$. A tape splice
(\ui{Klebestelle}) at every loop return adds a tick of noise decaying as $e^{-j/(0.0006f_s)}$ and a dip of depth
$0.85g$ that recovers linearly over \SI{12}{ms}.

\paragraph{Into the output.} The finished block $x_n$ ($n = 0\dots N-1$) is added to every output channel $c$ with a
gain that moves linearly from the previous block's value $g_c^0$ to this block's $g_c^1$ (the encoder gains of
\S\ref{sec:math-encoder} times all level factors). Since
\begin{equation}
  o_c[n] \mathrel{+}= x_n\bigl(g_c^0 + \Delta g_c\,n\bigr) = g_c^0\,x_n + \Delta g_c\,(n\,x_n),
\end{equation}
the product $n\,x_n$ is computed once per voice and the update of each of up to 36 channels is two vectorised
multiply-adds. The ramp removes the zipper noise that switching gains every 64 samples would cause, and it is why a
pipe can circle at \SI{500}{Hz}: between control points its direction is interpolated. For \ui{Sprünge} set to hard the
gains are held constant and switched at the exact sample of each jump, so the jumps click as intended.

\paragraph{The gains of one voice (\fn{updateVoiceSpatial}).} Once per control block the voice computes its position
(layout, motion, modulation offsets $\varphi \mathrel{+}= 180^\circ m_\varphi$, $\vartheta \mathrel{+}= 90^\circ m_\vartheta$,
$d \mathrel{*}= 2^{3m_d}$), moves it by the walk of the listener, and from it the gain
\begin{equation}
  g = g_\mathrm{stop}\cdot 10^{36m_L/20}\cdot g_d(d)\cdot 10^{-h/20}\cdot\cos\bigl(\tfrac\pi2 c_\mathrm{mix}\bigr),
\end{equation}
the product of the stop's level, the level modulation ($\pm\SI{36}{dB}$ per unit), the distance law, the loss $h$ in
dB when the pipe is hidden behind a corner of the floor plan (\S\ref{sec:plans}), and the equal-power share left to
the dry pipe when part of it dissolves into the grain cloud. The send to the room uses the much flatter law $g_d$ with
the exponent $0.35a$: the reverberant field of a room hardly depends on the distance to the source. Behind a corner
the direct sound also passes an extra low-pass whose coefficient glides over \SI{30}{ms},
$a \leftarrow a + (a^* - a)(1 - e^{-N/(0.03f_s)})$, because switching it at once would make a clear signal jump to a
dull one.

\paragraph{Doppler.} The radial speed is the change of distance per control block,
$v_r = (d_n - d_{n-1})/T_b$ (clamped to $\pm\SI{120}{m/s}$), and the reading rate is multiplied by
\begin{equation}
  D^* = \operatorname{clamp}_{[0.5,2]}\Bigl(\frac{c}{c + \eta\,v_r}\Bigr),\qquad c = \SI{343}{m/s},
\end{equation}
the frequency a listener at rest hears from a moving source, scaled by the amount $\eta$ (\ui{Doppler}). $D$ glides
towards $D^*$ with $a = 0.08$ per 16 samples ($\tau\approx\SI{4}{ms}$), so the steps of the block-wise derivative are
not heard. Jumping pipes get no Doppler, since their distance changes without travelling.

\paragraph{A pipe at the ear (\ui{Flüstern}).} When a pipe comes within $\SI{1.4}{m}$ of the head, it is also sent
straight to the two ears, which sit $\pm\SI{8.75}{cm}$ to the sides of the head turned by the view's yaw. With the
distances $d_L, d_R$ to the ears, $d_\mathrm{min} = \min(d_L,d_R)$, and the weight
$t = \operatorname{smoothstep}\bigl((1.4-d)/1.2\bigr)$ (the smoothstep $3t^2-2t^3$ starts and ends with zero slope):
\begin{align}
  g_{L,R} &= 0.9\,W\,t\,g_\mathrm{stop}\cdot\min\bigl(12,\ 1/d_{L,R}\bigr)\cdot(d_\mathrm{min}/d_{L,R})^{0.7}, \\
  \tau_{L,R} &= (d_{L,R} - d_\mathrm{min})\,f_s/c \quad\text{samples},\\
  f_{c,L,R} &= \operatorname{clamp}_{[1.5\mathrm{k},\,18\mathrm{k}]}\Bigl(\frac{\SI{18}{kHz}}{1 + 4(d_{L,R}/d_\mathrm{min} - 1)}\Bigr).
\end{align}
The $1/d$ law per ear makes the level difference between the ears large, as only a very near source produces; the
extra factor $(d_\mathrm{min}/d)^{0.7}$ is the head's shadow on the far ear, which at this distance reaches down to
low frequencies; the delay is the true difference of travel time, read with linear interpolation from a 128-sample
history; and the far ear is darkened by a one-pole whose corner falls with the ratio of the distances. A
\SI{180}{Hz} one-pole added with the weight $0.9t$ gives the bass boost of a source at the ear. Meanwhile the pipe's
normal gain is reduced by $1-0.85Wt$, so the two pictures do not double.

% ---------------------------------------------------------------------
\subsection{Degradation (\fn{Degrade.h})}
\label{sec:math-degrade}
\S\ref{sec:degrade} gives the formula of the chain; here are its parts as the code computes them.
\fn{DegradeParams::prepareBlock} derives once per block the quantiser step count $q = 2^{B-1}$ (off for
$B\ge 15.9$), the hold increment $h = \operatorname{clamp}(f_h, 1, f_s)/f_s$, the per-sample probability of a crackle
$40\gamma^2/f_s$ (on average $40\gamma^2$ per second) and of a dropout $3\delta^2/f_s$.
\fn{DegradeState::frame} then draws, once per sample frame and shared by all channels of the sum:
\begin{itemize}[nosep]
  \item \emph{sample and hold}: a phase accumulator $\phi \mathrel{+}= h$; when it passes 1 a new value is taken. The
    held value is not interpolated, so its images above $f_h/2$ remain: the metallic aliasing of a cheap converter is
    the point.
  \item \emph{crackle}: on an event an impulse of random sign and size $0.5$--$1.5$, then decaying by $0.93$ per sample
    ($\tau\approx14$ samples, a click of a fraction of a millisecond).
  \item \emph{dropout}: on an event a gap of \SIrange{20}{250}{ms}; the gain moves towards 0 and back to 1 by at most
    $1/(0.003f_s)$ per sample, i.e.\ linear \SI{3}{ms} ramps that do not click.
\end{itemize}
\fn{DegradeState::process} applies the frame to one sample: hold, round to the nearest of $2q$ levels in $[-1,1]$
($y = \lfloor yq + \tfrac12\rfloor/q$, rounding half away from zero), add the grain
$\gamma\,e\,(0.6n + 1.5c)$ where the envelope $e \mathrel{+}= 0.002(|x|-e)$ ($\tau = 500$ samples, about
\SI{10}{ms}) makes the noise follow the signal, multiply by the dropout gain, and mix
$x + \alpha(y-x)$. Because the grain is proportional to the envelope, silence stays silent: the decay lives on the
sound, like grain on a photograph lives on the exposed parts.
